\documentclass[11pt]{article}
% ===========================================================================
%  Shared preamble for the AP Calculus AB exam-preparation worksheets.
%  Input right after \documentclass. Files beginning with "_" are not compiled
%  on their own. Navy and gold, to match the course's lesson pages.
% ===========================================================================
\usepackage[letterpaper,top=2.3cm,bottom=2.2cm,left=1.7cm,right=1.7cm,headheight=28pt,footskip=30pt]{geometry}
\usepackage{amsmath,amssymb}
\usepackage[T1]{fontenc}
\usepackage[utf8]{inputenc}
\usepackage{lmodern}
\usepackage{xcolor}
\usepackage[most]{tcolorbox}
\usepackage{enumitem}
\usepackage{fancyhdr}
\usepackage{array}
\usepackage{tikz}
\usepackage{pgfplots}
\pgfplotsset{compat=1.18}

\definecolor{navy}{HTML}{161E54}
\definecolor{ink}{HTML}{27357E}
\definecolor{gold}{HTML}{B98A00}
\definecolor{goldbg}{HTML}{FFFDF3}
\definecolor{inkbg}{HTML}{F6F7FD}
\definecolor{rule}{HTML}{C9CFF0}
\definecolor{okgreen}{HTML}{1B7A44}
\definecolor{okbg}{HTML}{F1F8F3}

\pagestyle{fancy}
\fancyhf{}
\fancyhead[L]{\footnotesize Name: \underline{\hspace{4cm}}}
\fancyhead[C]{\textbf{\color{navy}Integration Academy}\\[-2pt]{\scriptsize Calculus AB \textperiodcentered{} AP\textsuperscript{\textregistered} Exam Preparation}}
\fancyfoot[L]{\tiny\color{black!55}\parbox[t]{0.78\textwidth}{AP\textsuperscript{\textregistered} is a registered trademark of the College Board, which was not involved in the production of, and does not endorse, this material. \textcopyright{} Integration Academy.}}
\fancyfoot[R]{\scriptsize\thepage}
\renewcommand{\headrulewidth}{1.2pt}
\renewcommand{\headrule}{\hbox to\headwidth{\color{gold}\leaders\hrule height \headrulewidth\hfill}}
\renewcommand{\footrulewidth}{0pt}
\setlength{\parindent}{0pt}
\setlength{\parskip}{4pt}

% Title block: \wstitle{code}{title}{unit line}
\newcommand{\wstitle}[3]{\begin{center}{\small\bfseries\color{gold}WORKSHEET #1}\\[3pt]{\LARGE\bfseries\color{navy}#2}\\[3pt]{\small\color{black!60}#3}\end{center}\vspace{2pt}}
% Section heading
\newcommand{\sech}[1]{\par\vspace{10pt}{\large\bfseries\color{navy}#1}\par\nobreak\vspace{2pt}{\color{rule}\hrule height0.8pt}\vspace{7pt}}
% A level tag printed beside a group of questions
\newcommand{\level}[1]{\par\vspace{6pt}{\footnotesize\bfseries\color{gold}\MakeUppercase{#1}}\par\nobreak\vspace{2pt}}

\newtcolorbox{ideabox}{colframe=ink,colback=inkbg,coltitle=white,colbacktitle=ink,fonttitle=\bfseries\small,title={Key ideas},arc=2pt,boxrule=0.8pt,left=8pt,right=8pt}
\newtcolorbox{example}[1]{colframe=rule,colback=white,coltitle=navy,colbacktitle=inkbg,fonttitle=\bfseries\small,title={#1},arc=2pt,boxrule=0.8pt,borderline west={3pt}{0pt}{ink},left=10pt,right=8pt}
\newtcolorbox{frbox}[1]{colframe=gold,colback=goldbg,coltitle=white,colbacktitle=gold,fonttitle=\bfseries\small,title={#1},arc=2pt,boxrule=0.8pt,breakable,left=8pt,right=8pt}
\newtcolorbox{solbox}{colframe=okgreen,colback=okbg,coltitle=white,colbacktitle=okgreen,fonttitle=\bfseries\small,title={Worked solutions},arc=2pt,boxrule=0.8pt,breakable,left=8pt,right=8pt}
\newtcolorbox{scorebox}[1]{colframe=navy,colback=white,coltitle=white,colbacktitle=navy,fonttitle=\bfseries\small,title={#1},arc=2pt,boxrule=0.8pt,breakable,left=8pt,right=8pt}

\newcommand{\soln}{\par\smallskip\textbf{\color{ink}Solution.}\ }
\newcommand{\result}[1]{\par\smallskip\textbf{#1}}
% blank working space of a given height
\newcommand{\work}[1]{\par\vspace{#1}}
% a point value in the margin of a scoring guide
\newcommand{\pts}[1]{\hfill{\footnotesize\bfseries\color{gold}[#1]}}

% Coordinate grid for graph questions: \begin{apgraph}{xmin}{xmax}{ymin}{ymax} ... \end{apgraph}
\newenvironment{apgraph}[4]{\begin{center}\begin{tikzpicture}\begin{axis}[
  axis lines=middle, xmin=#1, xmax=#2, ymin=#3, ymax=#4,
  grid=both, grid style={black!12}, xtick distance=1, ytick distance=1,
  ticklabel style={font=\scriptsize}, xlabel={$x$}, ylabel={$y$},
  width=10cm, height=7.2cm, samples=80, clip=false,
  every axis plot/.append style={line width=1.1pt, color=ink}]}{\end{axis}\end{tikzpicture}\end{center}}
\newcommand{\closeddot}[2]{\addplot[only marks, mark=*, mark size=2.2pt, color=ink] coordinates {(#1,#2)};}
\newcommand{\opendot}[2]{\addplot[only marks, mark=*, mark size=2.2pt, color=ink, mark options={fill=white}] coordinates {(#1,#2)};}

% Multiple-choice options: four across, or stacked when the options are long.
\newcommand{\opts}[4]{\par\vspace{3pt}\begin{tabular}{@{}p{0.235\linewidth}p{0.235\linewidth}p{0.235\linewidth}p{0.235\linewidth}@{}}(A)\ #1 & (B)\ #2 & (C)\ #3 & (D)\ #4\end{tabular}\par\vspace{9pt}}
\newcommand{\optsv}[4]{\par\vspace{3pt}\begin{tabular}{@{}l@{\ \ }p{0.9\linewidth}@{}}(A) & #1\\(B) & #2\\(C) & #3\\(D) & #4\end{tabular}\par\vspace{9pt}}

\fancyhead[R]{\footnotesize\color{black!60}Worksheet 2.3}
\begin{document}

\wstitle{2.3}{Estimating a Derivative from Data and Graphs}{Unit 2: The Derivative}

\begin{ideabox}
\begin{itemize}[leftmargin=*,itemsep=2pt,topsep=2pt]
  \item From a table, estimate $f'(c)$ with the slope between the \textbf{nearest values on either side} of $c$. Show the difference quotient.
  \item From a graph, $f'(c)$ is the slope of the graph at $c$. At a corner the derivative does not exist.
  \item An estimate of a derivative carries units, and can be interpreted like any other rate.
  \item A known value and a known rate give a linear estimate nearby: $f(c+h)\approx f(c)+f'(c)\,h$.
\end{itemize}
\end{ideabox}

\sech{Worked examples}
\begin{example}{Example 1 \textperiodcentered{} From a table}
The distance travelled by a train, $D(t)$ km, is recorded at times $t$ minutes.\begin{center}\renewcommand{\arraystretch}{1.25}\begin{tabular}{|c|c|c|c|c|}\hline $t$ & $0$ & $5$ & $10$ & $20$ \\ \hline $D(t)$ & $0$ & $6$ & $15$ & $40$ \\ \hline\end{tabular}\end{center}Estimate $D'(7.5)$ with units.\soln Use $t=5$ and $t=10$: $\dfrac{15-6}{10-5}=1.8$.\result{$D'(7.5)\approx1.8$ km per minute.}
\end{example}

\begin{example}{Example 2 \textperiodcentered{} A symmetric estimate}
Given $f(2.9)=7.1$, $f(3)=7.5$ and $f(3.1)=8.1$, estimate $f'(3)$.\soln The values on either side of 3 give $\dfrac{8.1-7.1}{3.1-2.9}$.\result{$f'(3)\approx5$.}
\end{example}

\begin{example}{Example 3 \textperiodcentered{} From a graph of line segments}
The graph of $g$ contains the line segment from $(1,2)$ to $(4,8)$. Find $g'(2)$.\soln On a line segment the derivative is the slope of the segment: $\dfrac{8-2}{4-1}$.\result{$g'(2)=2$.}
\end{example}

\newpage
\sech{Practice}
Questions 1 to 7 use this table. The population of a town, $P(t)$ people, is a differentiable function of time $t$ years.\begin{center}\renewcommand{\arraystretch}{1.25}\begin{tabular}{|c|c|c|c|c|c|}\hline $t$ & $0$ & $2$ & $4$ & $7$ & $10$ \\ \hline $P(t)$ & $100$ & $112$ & $130$ & $160$ & $172$ \\ \hline\end{tabular}\end{center}
\level{Core}
\begin{enumerate}[leftmargin=*,itemsep=2pt]
  \item Estimate $P'(3)$.\work{2.3cm}
  \item Estimate $P'(8.5)$.\work{2.3cm}
  \item Find the average rate of change of $P$ on $[0,10]$.\work{2.3cm}
  \item State the units of $P'(t)$.\work{2.3cm}
\end{enumerate}

\level{Exam level}
\begin{enumerate}[leftmargin=*,itemsep=2pt,start=5]
  \item Estimate $P'(5.5)$ and interpret your answer in context.\work{2.8cm}
  \item Use $P(4)=130$ and your estimate from Question 5 to approximate $P(4.5)$.\work{2.8cm}
  \item Find the average rate of change of $P$ on each of the four subintervals. What do the values suggest about how the growth of the town is changing?\work{2.8cm}
  \item The graph of $f$ consists of the line segments joining $(0,1)$, $(2,5)$ and $(5,-1)$. Find $f'(1)$ and $f'(3)$, and explain why $f'(2)$ does not exist.\work{2.8cm}
\end{enumerate}

\begin{frbox}{Question 9 \textperiodcentered{} Free-response practice}
The temperature along a metal rod, $T(x)$ degrees Celsius, is a differentiable function of the distance $x$ centimetres from one end.\begin{center}\renewcommand{\arraystretch}{1.25}\begin{tabular}{|c|c|c|c|c|c|}\hline $x$ & $0$ & $2$ & $5$ & $9$ & $12$ \\ \hline $T(x)$ & $80$ & $71$ & $62$ & $54$ & $51$ \\ \hline\end{tabular}\end{center}
\begin{enumerate}[label=(\alph*),leftmargin=*,itemsep=2pt]
  \item Estimate $T'(7)$. Indicate units of measure.\work{2.6cm}
  \item Interpret your answer to part (a) in the context of the problem.\work{2.6cm}
  \item Must there be a point of the rod at which $T'(x)=-\dfrac{29}{12}$? Justify your answer.\work{2.6cm}
\end{enumerate}
\end{frbox}

\newpage
\begin{solbox}
\begin{enumerate}[leftmargin=*,itemsep=4pt]
  \item $\dfrac{130-112}{4-2}=9$ people per year.
  \item $\dfrac{172-160}{10-7}=4$ people per year.
  \item $\dfrac{172-100}{10}=7.2$ people per year.
  \item People per year.
  \item $\dfrac{160-130}{7-4}=10$. At $t=5.5$ years the population is increasing at about 10 people per year.
  \item $130+10(0.5)=135$ people.
  \item $6$, $9$, $10$ and $4$ people per year. The growth speeds up and then slows down, which suggests the graph of $P$ changes from concave up to concave down.
  \item $f'(1)=\dfrac{5-1}{2}=2$ and $f'(3)=\dfrac{-1-5}{3}=-2$. At $x=2$ the slopes from the left and right are 2 and $-2$: a corner.
  \item \begin{enumerate}[label=(\alph*),leftmargin=*,itemsep=3pt]
    \item $\dfrac{54-62}{9-5}=-2$ degrees Celsius per centimetre.
    \item At the point 7 cm from the end, the temperature is decreasing at about 2 degrees Celsius per centimetre of distance along the rod.
    \item Yes. $\dfrac{T(12)-T(0)}{12}=-\dfrac{29}{12}$, and $T$ is differentiable (so continuous) on $[0,12]$, so the Mean Value Theorem guarantees such a point.
  \end{enumerate}
\end{enumerate}
\end{solbox}

\end{document}
